\documentclass[lettersize,journal]{IEEEtran}
\usepackage{amsmath,amsfonts}
\usepackage{algorithmic}
\usepackage{algorithm}
\usepackage{array}
\usepackage[caption=false,font=normalsize,labelfont=sf,textfont=sf]{subfig}
\usepackage{textcomp}
\usepackage{stfloats}
\usepackage{url}
\usepackage{verbatim}
\usepackage{graphicx}
\usepackage{cite}
\hyphenation{op-tical net-works semi-conduc-tor IEEE-Xplore}
% updated with editorial comments 8/9/2021
\usepackage{hyperref}
\usepackage{tabularx}
\usepackage{booktabs}

\begin{document}

\title{Order-aware Self-Cross Attention Representation Learning for Multi-omic Cancer Subtype Classification}

% \author{Tuan Duc Nguyen$^\dagger$, 
%         Hai Duc Tran$^\dagger$, 
%         and Duc-Hau Le$^*$%
% \thanks{$^\dagger$Tuan Duc Nguyen and Hai Duc Tran contributed equally to this work.}%
% \thanks{$^*$Corresponding author: \href{mailto:hauld@soict.hust.edu.vn}{hauld@soict.hust.edu.vn}}%
% \thanks{The authors are with the School of Information and Communications Technology, Hanoi University of Science and Technology, 1 Dai Co Viet, Bach Mai, Hanoi 100000, Vietnam.}}


\author{Tuan Duc Nguyen$^\dagger$, 
        Hai Duc Tran$^\dagger$, 
        and Duc-Hau Le$^*$%
\thanks{$^\dagger$Tuan Duc Nguyen and Hai Duc Tran contributed equally to this work.}% <-this % stops a space
\thanks{$^*$Corresponding author: \href{mailto:hauld@soict.hust.edu.vn}{hauld@soict.hust.edu.vn}}%
\thanks{The authors are with the School of Information and Communications Technology, Hanoi University of Science and Technology, 1 Dai Co Viet, Bach Mai, Hanoi 100000, Vietnam.}}


% The paper headers
\markboth{Journal of \LaTeX\ Class Files,~Vol.~14, No.~8, August~2021}%
{Shell \MakeLowercase{\textit{et al.}}: A Sample Article Using IEEEtran.cls for IEEE Journals}

% \IEEEpubid{0000--0000/00\$00.00~\copyright~2021 IEEE}
% Remember, if you use this you must call \IEEEpubidadjcol in the second
% column for its text to clear the IEEEpubid mark.  

\maketitle

\begin{abstract}
The precise classification of cancer subtypes is a cornerstone of precision oncology, yet the integration of heterogeneous multi-omic data remains a significant computational challenge. Existing deep learning frameworks often utilize symmetric fusion strategies that treat omic modalities as independent, parallel inputs, thereby overlooking the directional regulatory dependencies inherent in biological systems. We propose \textbf{OSCAR} (\textbf{O}rder-aware \textbf{S}elf-\textbf{C}ross \textbf{A}ttention \textbf{R}epresentation), a novel hierarchical deep learning framework that explicitly incorporates the biological ``Central Dogma'' flow: $\text{Genomics} \rightarrow \text{Transcriptomics} \rightarrow \text{Proteomics}$. Our architecture first encodes each modality with a shared-dimension MLP, then applies an Ordered Cross-Attention (OCA) mechanism in which transcriptomic signals (mRNA) are contextualized by genomic structural variations (CNV), and proteomic representations (RPPA) are subsequently guided by transcriptomic context. To account for post-transcriptional dysregulation and biological noise, we integrate a Gated Fusion module that adaptively balances hierarchical context with modality-specific signals. We evaluated OSCAR on three benchmark cohorts from The Cancer Genome Atlas (TCGA): breast invasive carcinoma (BRCA), lung cancer (LUNG), and colorectal adenocarcinoma (COADREAD). Our model achieves state-of-the-art performance, reaching an accuracy of \textbf{91.27\%} in the high-heterogeneity BRCA dataset and \textbf{96.71\%} in LUNG subtype classification, outperforming established graph-based and symmetric attention baselines. Furthermore, five-fold consensus input-gradient saliency, supported by differential abundance and pathway enrichment, shows that OSCAR recovers established clinical programs---including luminal mitotic/proliferation signatures and the 17q12/\textit{ERBB2} amplicon in HER2-enriched breast cancer---while also nominating less canonical but biologically plausible candidates such as \textit{TBC1D31} and \textit{GAS2L3}. These results demonstrate that incorporating directional biological priors through ordered attention not only enhances classification accuracy but also provides a robust, interpretable tool for cross-omic biomarker prioritization and clinical diagnostic support.
\end{abstract}

\begin{IEEEkeywords}
Multi-omic integration, Cancer subtype classification, Ordered cross-attention, Central Dogma, Gated fusion, Precision oncology.
\end{IEEEkeywords}

% \section{Introduction}
% \IEEEPARstart{T}{his} file is intended to serve as a ``sample article file''
% for IEEE journal papers produced under \LaTeX\ using
% IEEEtran.cls version 1.8b and later. The most common elements are covered in the simplified and updated instructions in ``New\_IEEEtran\_how-to.pdf''. For less common elements you can refer back to the original ``IEEEtran\_HOWTO.pdf''. It is assumed that the reader has a basic working knowledge of \LaTeX. Those who are new to \LaTeX \ are encouraged to read Tobias Oetiker's ``The Not So Short Introduction to \LaTeX ,'' available at: \url{http://tug.ctan.org/info/lshort/english/lshort.pdf} which provides an overview of working with \LaTeX.

\section{Introduction}
\label{sec:introduction}
\IEEEPARstart{T}{he} TCGA Pan-Cancer analyses have established molecular cancer classification as a central component of precision oncology, supporting prognosis, treatment stratification, and biomarker discovery, as framed by The Cancer Genome Atlas (TCGA) \cite{TCGA2013} and the Pan-Cancer Atlas \cite{Hoadley2018PanCancer}. Cancer is inherently heterogeneous; tumors arising in the same organ can follow markedly different clinical trajectories because they are driven by distinct molecular programs. Although histopathology remains clinically important, high-throughput profiling has shifted subtype discovery toward molecular signatures derived from genomics, transcriptomics, epigenomics, and proteomics. Large-scale TCGA studies of breast, lung, and colorectal tumors have shown that molecular subtypes often refine or complement tissue-level diagnosis, including TCGA breast characterization \cite{Koboldt2012BRCA}, LUAD \cite{TCGA2014LUAD}, LUSC \cite{TCGA2012LUSC}, and COADREAD \cite{TCGA2012COADREAD}. 

Recent reviews \cite{Cai2022Review, CarrilloPerez2022, Abdelaziz2024Review, Han2025Review} have emphasized that multi-omic integration is now central to cancer research because single-omic views rarely capture the layered regulatory mechanisms that connect genotype to phenotype. Integrating complementary modalities such as Copy Number Variations (CNV), messenger RNA (mRNA), and Reverse Phase Protein Array (RPPA) is therefore essential for a more faithful representation of oncogenic processes and subtype-specific biology \cite{CrossAttOmics2024, Omran2025}. 

Over the past decade, the computational strategies for multi-omic fusion have evolved through distinct technological paradigms. Early methodologies primarily relied on matrix factorization and joint statistical modeling—such as iCluster \cite{iCluster2009}, integrative NMF \cite{IntNMF2017}, and JIVE \cite{JIVE2013}—or simple feature concatenation extended by deep networks like OmiEmbed \cite{OmiEmbed2021} and DeepMO \cite{DeepMO2020}. While providing a rigorous statistical foundation, these early methods frequently struggle with non-linear cross-modal dependencies and are highly vulnerable to scale mismatches and noise propagation. To inject structural awareness, the literature shifted toward graph-based architectures; models like MOGONET \cite{MOGONET2021}, MoGCN \cite{MoGCN2022}, DeepMoIC \cite{DeepMoIC2024}, and $M^2$CGCN \cite{M2CGCN2025} leverage patient similarity networks or feature graphs to refine inter-omic correlations. Concurrently, the success of Transformers inspired attention-based frameworks—including moBRCA-net \cite{moBRCAnat2023}, CrossAttOmics \cite{CrossAttOmics2024}, MoxGATE \cite{MoxGATE2023}, and the Multi-Omics Transformer (MOT) \cite{MOT2023}—which utilize symmetric cross-attention or gated mechanics to adaptively align heterogeneous features, while generative or pathway-guided networks \cite{SubtypeGAN2020, SurvPath2023} attempted to constrain these latent representations using biological priors.

Despite substantial methodological progress, a key conceptual limitation persists across these modern paradigms. Most existing models, whether graph-based, attention-driven, or generative, still treat multi-omic modalities as symmetric or parallel inputs. This ``flat" fusion assumption overlooks the directional regulatory logic summarized by the Central Dogma, in which genomic alterations influence transcript abundance and transcripts subsequently shape protein-level outcomes. In practice, CNVs, mRNA expression, and RPPA measurements are not merely correlated feature blocks; they are connected by sequential dependencies. Consequently, even state-of-the-art (SOTA) methods may miss subtype-discriminative signals that depend on cross-layer ordering rather than on undirected statistical similarity alone. Furthermore, this limitation propagates into model explainability. Traditional post-hoc tools (e.g., SHAP, LIME) or modern symmetric attention weights \cite{moBRCAnat2023, OmiEmbed2021} evaluate features either in isolation or as undirected correlations, failing to capture the directional, cascading nature of cross-omic interactions.

To address this gap, we propose OSCAR (Order-aware Self-Cross Attention Representation), a deep learning architecture that explicitly incorporates biological hierarchy through an Ordered Cross-Attention (OCA) mechanism. Unlike recent integration models that treat modalities as parallel inputs, OSCAR introduces a directional inductive bias aligned with the regulatory flow from CNV to mRNA to RPPA. Our framework first maps each omic profile to a latent vector, then performs hierarchical ordered cross-attention so that genomic structural variation provides context for transcriptomic interpretation, which in turn contextualizes proteomic representation, followed by inter-omic self-attention over the resulting modality tokens. This design grounds the interaction mechanism in a biologically plausible sequence that directly motivates the present study.

\begin{figure*}[!t]
    \centering
    \includegraphics[width=1\textwidth]{model.pdf}
    \caption{\textbf{Proposed OSCAR Architecture.} The model captures the biological hierarchy from Genomics to Transcriptomics and Transcriptomics to Proteomics through ordered cross-attention, gated fusion, and inter-omic self-attention.}
    \label{fig:model}
\end{figure*}

Furthermore, we introduce an adaptive Gated Fusion mechanism to balance hierarchical context with modality-specific signals, preventing the loss of unique biomarkers due to post-transcriptional noise. The primary contributions of this work are three-fold:
\begin{enumerate}
\item We propose OSCAR, an order-aware self-cross attention framework for multi-omic fusion that honors the directional flow of biological information.
\item We demonstrate that our model achieves better performance across three diverse TCGA cohorts: breast invasive carcinoma (BRCA), lung cancer (LUAD/LUSC, collectively denoted as LUNG), and colorectal adenocarcinoma (COAD/READ, denoted as COADREAD), particularly in resolving high-heterogeneity subtypes.
\item We provide a comprehensive interpretability analysis based on multi-omic input-gradient saliency with cross-fold consensus, differential abundance filtering, and pathway enrichment, showing alignment with established clinical biomarkers and nominating additional subtype-informative candidates.
\end{enumerate}



% 



\section{Materials and Methods}

\subsection{Data Preparation}
We evaluated OSCAR on three TCGA multi-omic cohorts: breast invasive carcinoma (BRCA), lung cancer (LUNG)---comprising lung adenocarcinoma (LUAD) and lung squamous cell carcinoma (LUSC)---and colorectal adenocarcinoma (COADREAD). Each cohort includes matched genomic copy number variation (CNV), transcriptomic mRNA expression, proteomic reverse phase protein array (RPPA), and clinical subtype annotations. These cohorts were selected because their molecular landscapes and subtype definitions have been extensively characterized by the TCGA consortium, including benchmark studies on BRCA \cite{Koboldt2012BRCA}, LUAD \cite{TCGA2014LUAD}, LUSC \cite{TCGA2012LUSC}, and COADREAD \cite{TCGA2012COADREAD}. Molecular and clinical data were obtained from the UCSC Xena platform \cite{Goldman2020Xena}, with proteomic profiles supplemented from The Cancer Proteome Atlas (TCPA) \cite{Li2013TCPA}.

During preprocessing, we retained only patients shared across CNV, mRNA, RPPA, and clinical metadata to ensure complete multi-omic alignment. Samples without subtype labels and features containing missing values were removed. The resulting cohorts represent a high-dimensional, small-sample classification setting, with more than 45,000 molecular features and fewer than 700 aligned patients per dataset, as summarized together with the subtype distributions in Table \ref{tab:dataset_summary}. 

\begin{table*}[htbp]
\centering
\caption{Summary of experimental datasets, multi-omic feature dimensions, and subtype distributions.}
\label{tab:dataset_summary}
\small
\begin{tabular}{@{}lccrX@{}}
\toprule
\textbf{Dataset} & \textbf{Modality} & \textbf{Features ($D$)} & \textbf{Subtype distribution} \\ 
\midrule
\textbf{LUNG} (N=682) & mRNA & 20,530 & LUAD: 359 (52.6\%) \\
 &  CNV & 24,776 & LUSC: 323 (47.4\%) \\
 &  RPPA & 180 &  \\ 
\midrule
% \textbf{BRCA} (N=511) & mRNA & 20,530 & Luminal A: 248 (49.0\%) \\
%  & CNV & 24,776 & Luminal B: 98 (24.4\%) \\
%  & RPPA & 168 &  Basal-like: 112 (17.4\%) \\ 

%  &  &  & HER2-enriched: 53 (9.1\%) \\ 
% %  &  &  & Normal-like: 16 (2.3\%) \\ 

\textbf{BRCA} (N=511) & mRNA & 20,530 & Luminal A: 248 (48.5\%) \\
 & CNV & 24,776 & Luminal B: 98 (19.2\%) \\
 & RPPA & 168 & Basal-like: 112 (21.9\%) \\
 & & & HER2-enriched: 53 (10.4\%) \\
 
\midrule
\textbf{COADREAD} (N=222) & mRNA & 20,530 & CMS2: 86 (38.7\%)\\
 & CNV & 24,776 & CMS4: 69 (31.1\%)\\
 & RPPA & 153 & CMS3: 34 (15.3\%) \\ 
 &  &  & CMS1: 33 (14.9\%) \\ 

\bottomrule
\end{tabular}
\end{table*}

\subsection{Experimental Setup}
For model evaluation, we applied stratified 5-fold cross-validation to each cohort so that subtype proportions were preserved across folds. In each iteration, one fold was held out for testing and the remaining four folds were used for training. Reported performance values are averaged across the five runs, reducing dependence on any single train-test split. To evaluate the model’s performance, we report Accuracy (Acc) and F1 Score (F1).

Each modality was encoded into a shared latent space of $d=128$ via a ReLU MLP with a 512-dimensional hidden layer. Self-attention over omic tokens used 8 heads ($d_k=16$) and 2 Transformer layers. The classifier was a ReLU MLP with a 128-unit hidden layer and dropout 0.4. Models were trained on an NVIDIA Tesla V100 GPU for up to 200 epochs using AdamW, learning rate $2 \times 10^{-4}$, weight decay $1 \times 10^{-5}$, batch size 32 and Cross-Entropy loss.

The OSCAR framework, visualized in Figure \ref{fig:model}, integrates multi-omic data in four stages: per-omic MLP encoding, Ordered Cross-Attention (OCA), Gated Fusion, and inter-omic self-attention followed by subtype classification.

\subsubsection{Interpretability Analysis}
For biomarker prioritization on BRCA, we computed input-gradient saliency on held-out validation folds using the fold-wise best checkpoint.
Specifically, we attributed the true-class logit by the mean absolute gradient $\mathrm{Mean}\,|\partial\mathrm{logit}/\partial x|$ with respect to each CNV, mRNA, and RPPA input feature under evaluation mode.
Features that appeared in the per-omic top-$k$ list in at least three of five folds were retained as consensus candidates.
We further assessed one-versus-rest differential abundance (Mann--Whitney $U$ with Benjamini--Hochberg FDR and Cohen's $d$) and performed over-representation analysis with Enrichr using GO Biological Process, KEGG, and MSigDB Hallmark libraries.

\subsection{Problem Formulation}
Let the input space be defined by three distinct omic modalities from the same patient: genomic data (CNV) $X_{G} \in \mathbb{R}^{d_G}$, transcriptomic data (mRNA) $X_{T} \in \mathbb{R}^{d_T}$, and proteomic data (RPPA) $X_{P} \in \mathbb{R}^{d_P}$. Given a training set of $N$ samples, our goal is to learn a mapping $f: \{X_G, X_T, X_P\} \rightarrow \mathcal{Y}$, where $\mathcal{Y}$ corresponds to the ground-truth cancer subtype labels.

\subsection{Per-omic Latent Encoding}
Before cross-modal interaction, each high-dimensional omic profile is projected into a shared latent space. For modality $i \in \{G, T, P\}$, a two-layer ReLU MLP yields a single patient-level embedding
\begin{equation}
    Z_i = \mathrm{MLP}_i(X_i) \in \mathbb{R}^{d},
\end{equation}
with $d=128$. Because each modality is summarized as one vector rather than a variable-length token sequence, subsequent cross-modal interaction cannot rely on sequence Softmax attention without degenerating to a trivial weight of one; OSCAR therefore uses an element-wise ordered gating mechanism (OCA), and reserves multi-head self-attention for fusion over the three modality tokens after OCA.

\subsection{Ordered Cross-Attention (OCA)}
The defining innovation of OSCAR is the Ordered Cross-Attention (OCA) layer. Unlike recent symmetric integration models such as MoGCN \cite{MoGCN2022} or MoxGATE \cite{MoxGATE2023}, which treat modalities as parallel inputs, OSCAR models the directional hierarchy of the biological ``Central Dogma.''
Given a query embedding $Q$ and a key/value embedding from an upstream modality $K{=}V$, OCA applies learned projections, an element-wise sigmoid gate on the concatenated pair, and a residual update
\begin{equation}
    \tilde{Q}=W_Q Q,\quad \tilde{K}=W_K K,\quad \tilde{V}=W_V V,
\end{equation}
\begin{equation}
    \alpha = \sigma\!\left(\mathrm{MLP}\big([\tilde{Q}\parallel\tilde{K}]\big)\right),\qquad
    H = \tilde{Q} + \alpha \odot \tilde{V},
\end{equation}
followed by a position-wise feed-forward block with residual connections. The gate $\alpha \in (0,1)^{d}$ controls, dimension-wise, how much upstream (key/value) information is injected into the downstream query modality.

\subsubsection{Genomic-to-Transcriptomic Interaction}
In the first stage, transcriptomic features query genomic features so that CNV context modulates mRNA representation:
\begin{equation}
    H_T = \mathrm{OCA}(Q=Z_T,\ K{=}V=Z_G).
\end{equation}
This produces an enhanced transcriptomic embedding $H_T$ that is conditioned on the genomic blueprint, consistent with the regulatory direction $\mathrm{CNV}\rightarrow\mathrm{mRNA}$.

\subsubsection{Transcriptomic-to-Proteomic Interaction}
Following the hierarchy, proteomic features then query the (original) transcriptomic embedding:
\begin{equation}
    H_P = \mathrm{OCA}(Q=Z_P,\ K{=}V=Z_T).
\end{equation}
Using $Z_T$ rather than $H_T$ as the key/value source keeps the second stage anchored to the native transcriptomic signal while still enforcing the ordered dependency $\mathrm{mRNA}\rightarrow\mathrm{RPPA}$.


\begin{table*}[htbp]
\centering
\small
\caption{Classification performance comparison against existing methods (mean $\pm$ std \%). Bold indicates best performance.}
\label{tab:results}
\begin{tabular}{l|cc|cc|cc}
\hline
\textbf{Method} & \multicolumn{2}{c|}{\textbf{BRCA}} & \multicolumn{2}{c|}{\textbf{LUNG}} & \multicolumn{2}{c}{\textbf{COADREAD}} \\
 & F1 Score (\%) & Acc (\%) & F1 Score (\%) & Acc (\%) & F1 Score (\%) & Acc (\%) \\ \hline
MoGCN \cite{MoGCN2022} & 87.42 $\pm$ 2.15 & 86.91 $\pm$ 3.10 & 94.12 $\pm$ 1.88 & 93.85 $\pm$ 1.95 & 90.12 $\pm$ 2.45 & 89.34 $\pm$ 3.10 \\
moBRCA-net \cite{moBRCAnat2023} & 89.15 $\pm$ 1.20 & 88.54 $\pm$ 1.45 & -- & -- & -- & -- \\
MoxGATE \cite{MoxGATE2023} & 90.24 $\pm$ 3.01 & 89.73 $\pm$ 3.42 & 96.01 $\pm$ 2.05 & 95.87 $\pm$ 2.11 & 92.16 $\pm$ 1.22 & 91.67 $\pm$ 4.27 \\ 
MOGONET \cite{MOGONET2021} & 72.47 $\pm$ 2.60 & 71.10 $\pm$ 2.27 & 91.14 $\pm$ 3.18 & 91.45 $\pm$ 3.03 & 74.28 $\pm$ 5.27 & 73.28 $\pm$ 5.71 \\ \hline
\textbf{OSCAR} & \textbf{91.36 $\pm$ 2.26} & \textbf{91.27 $\pm$ 2.22} & \textbf{96.76 $\pm$ 1.37} & \textbf{96.71 $\pm$ 1.42} & \textbf{93.14 $\pm$ 3.28} & \textbf{92.34 $\pm$ 3.98} \\ \hline
\end{tabular}
\end{table*}

\subsection{Gated Fusion, Inter-omic Self-Attention, and Classification}
While the OCA layers provide hierarchical context, we must prevent the potential loss of unique modality-specific signals. For the downstream modalities $i \in \{T, P\}$, we implement a Gated Fusion mechanism that adaptively balances the OCA-enhanced embedding $H_i$ with the original latent $Z_i$:
\begin{equation}
    g_i = \sigma\!\left(\mathrm{MLP}_g\big([Z_i \parallel H_i]\big)\right),
\end{equation}
\begin{equation}
    F_i = g_i \odot H_i + (1 - g_i) \odot Z_i,
\end{equation}
where $\parallel$ denotes concatenation and $\odot$ denotes element-wise multiplication. The genomic embedding is kept as $F_G = Z_G$. The gate $g_i$ allows the network to decide, on a sample-by-sample basis, how much hierarchical information is required for accurate classification.

The fused modality tokens $\{F_G, F_T, F_P\}$ are then stacked and refined by a multi-head Transformer encoder~\cite{Vaswani2017Attention} (8 heads, 2 layers). A learned pooling (CLS) token aggregates the three omic tokens into a patient embedding $\mathcal{F} \in \mathbb{R}^{d}$, which is passed through a ReLU MLP classifier:
\begin{equation}
    \hat{y} = \text{Softmax}(\text{MLP}(\mathcal{F})).
\end{equation}
The model is optimized using the Cross-Entropy loss function across all $C$ classes.




\section{Results}
\subsection{Performance Evaluation and Comparison}

We present the main comparative results in Table \ref{tab:results}. Overall, the proposed OSCAR framework with Gated Fusion delivers the strongest and most consistent performance across the three benchmark cohorts. Importantly, OSCAR is not only competitive on a single preferred metric, but remains robust across both F1 score and accuracy, indicating that the model balances class-sensitive performance with global predictive reliability.



Several findings deserve emphasis. First, the results on BRCA provide the clearest evidence of the value of hierarchical multi-omic modeling. BRCA is the most heterogeneous cohort in our study, comprising four molecularly distinct subtypes with overlapping biological signatures. In this challenging setting, our model achieves the best overall performance, reaching 91.36\% F1 and 91.27\% accuracy. Compared with the specialized moBRCA-net \cite{moBRCAnat2023}, OSCAR improves the F1 score by 2.21\% and the accuracy by 2.73\%. It also surpasses the more general graph-based baseline MoGCN \cite{MoGCN2022} by even larger margins. These gains are important because BRCA is precisely the kind of dataset in which subtle cross-omic dependencies matter most; the improvement suggests that explicitly respecting the biological progression from genomics to transcriptomics and proteomics helps the model discriminate among closely related subtypes more effectively than flat fusion strategies.

Second, the method demonstrates strong generalization across different cancer types. On the LUNG cohort, OSCAR achieves 96.76\% F1 and 96.71\% accuracy, outperforming all competing methods on both metrics. This result is particularly notable because it shows that the proposed directional inductive bias does not overfit to one tissue context or one subtype taxonomy. Instead, it appears to capture a more general integration principle that remains effective even when the molecular landscape changes. The relatively small standard deviations on LUNG further indicate that the gains are stable across repeated runs rather than being driven by a favorable split or isolated experimental fluctuation.

Third, the COADREAD results reveal the superior performance of our proposed model across both class-balanced discrimination and global correctness. On this cohort, OSCAR achieves a higher F1 score (93.14\% versus 92.16\% for MoxGATE \cite{MoxGATE2023}), outperforming it by 0.98\%. Similarly, our model secures the best overall accuracy at 92.34\%, slightly exceeding MoxGATE's 91.67\%. This consistent improvement demonstrates that OSCAR is not only highly effective at capturing minority-subtype sensitivity (as reflected by the F1 score) but also more dependable in maintaining overall correctness across the entire patient cohort—a crucial requirement for practical clinical decision support systems.

Taken together, these results indicate that the main strength of OSCAR is not merely peak performance on a single dataset, but consistent superiority across heterogeneous settings. The proposed framework performs best in the most difficult cohort, remains dominant in the highest-accuracy cohort, and establishes a clear advantage even against highly competitive, specialized baselines. This overall pattern supports our central claim: incorporating directional biological structure into cross-omic attention provides a more reliable and transferable foundation for multi-omic cancer subtype classification than symmetric or purely similarity-based fusion mechanisms.

\subsection{Ablation Study}


To rigorously evaluate the contribution of each module within our framework, we conducted a series of ablation experiments on the BRCA and LUNG datasets. These experiments focus on two aspects: the architectural components and the hierarchical sequence of the attention mechanism.

\subsubsection{Importance of Model Components}
We compared the full Gated Fusion model against three architectural variants: the simple concatenation baseline (Concat), a version removing the inter-omic self-attention over modality tokens (w/o Self-Attention), and a version replacing the adaptive gate with a linear projection (w/o Gated Fusion). The results, summarized in Table \ref{tab:ablation_components}, illustrate the critical role of each module in achieving high classification accuracy.


\begin{table*}[!htbp]
\centering
\small
\caption{Ablation study of proposed model components (\% Accuracy).}
\label{tab:ablation_components}
\begin{tabular}{lccc}
\hline
\textbf{Model Variant} & \textbf{BRCA} & \textbf{LUNG} & \textbf{COADREAD} \\ \hline
Concat (baseline)      & 85.85 $\pm$ 4.08 & 93.43 $\pm$ 1.15 & 89.60 $\pm$ 3.60 \\
w/o Self-Attention     & 87.12 $\pm$ 2.45 & 94.88 $\pm$ 1.62 & 90.45 $\pm$ 2.80 \\
w/o Gated Fusion       & 88.24 $\pm$ 3.01 & 96.10 $\pm$ 1.31 & 91.62 $\pm$ 4.55 \\ \hline
\textbf{OSCAR (Full)} & \textbf{91.27 $\pm$ 2.22} & \textbf{96.71 $\pm$ 1.42} & \textbf{92.34 $\pm$ 3.98} \\ \hline
\end{tabular}
\end{table*}

The performance trends observed in Table \ref{tab:ablation_components} provide several key insights into multi-omic integration. The significant jump from the Concat baseline to the version utilizing inter-omic self-attention suggests that allowing modality tokens to exchange information after ordered cross-attention is important for effective fusion. By refining the gated CNV/mRNA/RPPA embeddings jointly, the model consolidates complementary signals before classification. 

The most substantial gains, however, are attributed to the Ordered Cross-Attention (OCA) and Gated Fusion mechanisms within OSCAR. Removing the Gated Fusion layer results in a 2.03\% accuracy drop in the BRCA dataset. This indicates that while the biological hierarchy provides a powerful structural prior, the model requires the flexibility to bypass this flow when the regulatory coupling between transcriptomics and proteomics is weakened—a common occurrence in aggressive cancer subtypes characterized by post-transcriptional dysregulation.

\subsubsection{Impact of Cross-Attention Order}
A unique hypothesis of our framework is that following the biological ``Central Dogma" sequence provides a superior inductive bias. To validate this, we compared our $Genomics\ (CNV) \rightarrow Transcriptomics\ (mRNA) \rightarrow Proteomics\ (RPPA)$ flow against alternative permutations. As shown in Table \ref{tab:ablation_order}, the biological order consistently outperforms mixed or reverse sequences across all datasets.

\begin{table*}[!htbp]
\centering
\small
\caption{Impact of different cross-attention orders on classification accuracy (\%).}
\label{tab:ablation_order}
\begin{tabular}{lccc}
\hline

\textbf{Ordering Strategy} & \textbf{BRCA} & \textbf{LUNG} & \textbf{COADREAD} \\ \hline
$Proteomics \rightarrow Transcriptomics \rightarrow Genomics$ & 86.44 $\pm$ 2.10 & 94.05 $\pm$ 1.85 & 89.09 $\pm$ 4.16 \\
(Reverse) \\
$Transcriptomics \rightarrow Genomics \rightarrow Proteomics $ & 88.27 $\pm$ 4.11 & 95.12 $\pm$ 1.40 & 90.43 $\pm$ 3.85 \\
(Mixed) \\
\textbf{$ Genomics \rightarrow Transcriptomics \rightarrow Proteomics $} & \textbf{91.27 $\pm$ 2.22} & \textbf{96.71 $\pm$ 1.42} & \textbf{92.34 $\pm$ 3.98} \\ 
\textbf{(OSCAR)}\\ 

\hline
\end{tabular}
\end{table*}


The failure of the ``Reverse" order (Proteomics (RPPA) $\rightarrow$ Transcriptomics (mRNA) $\rightarrow$ Genomics (CNV)) is particularly telling. Reversing which modality acts as query versus key/value is mathematically valid but biologically contradictory, leading to a loss of approximately 3.83\% in BRCA accuracy. This confirms that the directional flow of information is not merely a theoretical preference but a functional requirement for modeling multi-omic dependencies.

% \section{Model Explainability}
% To ensure clinical relevance, we analyzed the interpretability of OSCAR through attention weight visualization.

% \subsection{Feature Importance and Biomarker Discovery}
% By extracting the attention scores from the final layer, we identified the top-contributing features for each cancer subtype. For example, in the BRCA dataset, the model assigned high weights to the $ERBB2$ gene at the RNA level and its corresponding protein expression in HER2-enriched samples. This alignment with established clinical biomarkers validates that our model learns the underlying biology of cancer rather than mere statistical noise.

% \subsection{Cross-omic Correlation Analysis}
% We visualized the cross-attention maps between the $CNV$ and $RNA$ layers. The diagonal patterns observed in these maps suggest that the model successfully learns the "cis-acting" effects, where a copy number gain in a specific genomic region directly correlates with the over-expression of mRNA in that same region. These insights could potentially assist researchers in identifying novel regulatory drivers for rare cancer subtypes.


% \section{Discussion and Interpretability}

\subsection{Model Interpretability and Biological Implication}

The experimental results presented in Section IV demonstrate that the OSCAR framework consistently outperforms symmetric fusion models. In this section, we discuss the biological implications of our findings and provide an interpretability analysis to validate the clinical relevance of the learned features.

\subsubsection{Biological Rationale for Hierarchical Integration}
The primary innovation of our work is the implementation of a directional inductive bias that follows the flow of the biological central dogma. Most existing models, such as MoxGATE \cite{MoxGATE2023} and CrossAttOmics \cite{CrossAttOmics2024}, utilize symmetric cross-attention where modalities query each other in a bidirectional or ``flat" manner. While mathematically flexible, this approach ignores the causal relationship where genomic alterations (CNV) act as the template for transcriptomic (mRNA) activity, which in turn regulates protein synthesis (RPPA).

By letting transcriptomic embeddings query genomic embeddings (and proteomic embeddings subsequently query transcriptomic ones), our model injects upstream regulatory context into downstream representations. This is particularly relevant in cancer, where many mRNA fluctuations are stochastic ``noise" rather than oncogenic drivers. Our ablation study (Table \ref{tab:ablation_order}) confirmed that reversing this order significantly degrades performance, suggesting that the model benefits from a ``biological guide" that directs ordered cross-attention toward features with higher functional impact.

\subsubsection{Interpretability and Biomarker Discovery}

To move beyond a purely predictive ``black box,'' we quantified multi-omic feature importance with input-gradient saliency as described in Section~II.
Features retained by five-fold consensus were further filtered by differential abundance, and pathway coherence was assessed by Enrichr over-representation analysis.

\begin{figure*}[htbp]
    \centering
    \includegraphics[width=1\textwidth]{attention_plots/biomarker_step1_consensus/hbar_consensus_all_subtypes_annotated.png}
    \caption{Five-fold consensus gradient biomarkers (horizontal bars; color = omic).
    Soft background shading groups related genes (legend: mitotic / cell-cycle, 17q12 amplicon, signaling / lineage).
    Panels: \textbf{(A)} Luminal~A, \textbf{(B)} Luminal~B, \textbf{(C)} Basal-like, \textbf{(D)} HER2-enriched
    (top-12; panel-specific \emph{x}-scales)
    \cite{Gao2018LuminalA,Ades2014LuminalB,PerezPena2017MitoticLumA,Malorni2016RBsig,Staaf2010HER2,Kao2006Amplicon,Sahlberg2013HER2Amplicon}.}
    \label{fig:combined_saliency}
\end{figure*}

\paragraph{Recovery of established subtype programs.}
As summarized in Fig.~\ref{fig:combined_saliency}, OSCAR focuses on biologically recognized programs
\emph{without} injecting curated biomarker priors into training:

\begin{itemize}
    \item \textit{Luminal A and B.}
    Consensus features are enriched for mitotic apparatus and cell-cycle genes
    (e.g.\ \textit{ASPM}, \textit{CEP55}, \textit{CENPE}, \textit{NUF2}, \textit{DLGAP5}, \textit{CDKN3}, \textit{KIF11},
    \textit{GTSE1}, \textit{SGO1}), in line with proliferation-oriented distinctions among luminal tumors
    \cite{Gao2018LuminalA,Ades2014LuminalB,PerezPena2017MitoticLumA,Althobiti2024ASPM}.
    Hallmark over-representation further supports this reading, with significant enrichment for
    \textbf{E2F Targets}, \textbf{G2--M Checkpoint}, and \textbf{Mitotic Spindle}
    \cite{Malorni2016RBsig,Liberzon2015Hallmark}.
    Differential abundance clarifies polarity: mitotic transcripts are systematically \emph{lower} in Luminal~A versus
    the remainder and comparatively \emph{higher} in Luminal~B, matching the more proliferative luminal phenotype
    \cite{Ades2014LuminalB}.

    \item \textit{HER2-enriched.}
    Consensus features are dominated by the 17q12/\textit{ERBB2} amplicon neighborhood at both CNV and mRNA levels
    (e.g.\ \textit{STARD3}, \textit{TCAP}, \textit{MIEN1}, \textit{GRB7}, \textit{PGAP3}, \textit{PNMT}, \textit{CDK12}, \textit{GSDMB}),
    consistent with genomic characterization of HER2-positive disease
    \cite{Staaf2010HER2,Kao2006Amplicon,Sahlberg2013HER2Amplicon}.
    Multi-omic concordance for the same locus (CNV$\cap$mRNA) supports cis-acting amplification--expression coupling learned by the ordered hierarchy.

    \item \textit{Basal-like.}
    The basal panel is more mixed: gradient-stable features include transcripts such as \textit{MIA} and \textit{PGAP3},
    together with proteomic signals (e.g.\ GATA3 and phospho-TUBERIN/TSC2), rather than a single textbook amplicon driver.
    This heterogeneity is consistent with the known molecular diversity of basal-like / triple-negative disease
    \cite{Koboldt2012BRCA}, and we therefore emphasize basal attributions as model-facing prioritizations for follow-up
    rather than as standalone clinical biomarkers.
\end{itemize}

\paragraph{Model-nominated candidates beyond textbook markers.}
Beyond recovering canonical programs, the intersection of fold-stable gradients and differential abundance nominated
less stereotyped genes that are biologically plausible yet under-emphasized as subtype biomarkers:

\begin{itemize}
    \item \textit{TBC1D31} (mRNA) was consensus-stable for both Luminal~A (decreased vs.\ rest) and Luminal~B
    (increased vs.\ rest; present in $4/5$ folds for Luminal~B).
    TBC1-domain proteins regulate membrane trafficking / ciliogenesis and have recently been linked to metabolic
    plasticity in aggressive breast cancer, including a glycolytic role for \textit{TBC1D31} in TNBC models
    \cite{Arora2024TBC1D,Senatore2021TBC1D31}.
    OSCAR's unsupervised prioritization of \textit{TBC1D31} as a \emph{luminal A/B contrast feature} therefore
    extends this biology beyond a purely TNBC-centric narrative and nominates it for targeted follow-up.

    \item \textit{GAS2L3} (mRNA) was a strong Luminal~A consensus feature with large negative effect size versus
    non-Luminal~A samples.
    Functionally, GAS2L3 is a DREAM/E2F-regulated cytokinesis gene required for abscission and genomic stability
    \cite{Wolter2012GAS2L3,Peer2013GAS2L3}, providing a mechanistic fit to the Hallmark E2F / G2--M enrichment
    observed for Luminal~A, while remaining less commonly used than canonical Ki-67/Aurora-type proliferation markers
    in subtype biomarker lists.
\end{itemize}

\paragraph{Take-home.}
Collectively, (i)~alignment with established luminal proliferation and HER2-amplicon biology supports the
\emph{interpretability} of OSCAR~\cite{Staaf2010HER2,Gao2018LuminalA,Ades2014LuminalB}, and
(ii)~gradient$\cap$DE prioritization of \textit{TBC1D31} and \textit{GAS2L3} provides
\emph{hypothesis-generating} subtype-informative candidates for subsequent validation.
We do not claim these genes as new clinical biomarkers; rather, they illustrate how multi-omic attribution can surface
non-obvious but pathway-consistent features alongside known programs
\cite{Arora2024TBC1D,Wolter2012GAS2L3,PerezPena2017MitoticLumA}.

\subsubsection{The Role of Gated Fusion in Mitigating Noise}
A significant challenge in multi-omic integration is the inconsistent correlation between layers; for instance, high mRNA levels do not always lead to high protein levels due to post-transcriptional regulation. Our Gated Fusion mechanism addresses this by adaptively weighing the ``hierarchical" signal ($H_i$) against the ``raw" intra-omic signal ($Z_i$). 

Our analysis of the gate values $g_i$ revealed that for aggressive subtypes in the LUNG dataset, specifically \textit{Lung Squamous Cell Carcinoma (LUSC)}, the model relied more heavily on the raw mRNA signals than the hierarchical CNV-contextualized signals. This suggests that in advanced stages of lung squamous cell carcinoma, the regulatory link between copy number variations and gene expression becomes decoupled, and our model successfully adapts to this biological complexity by \textit{``bypassing''} the hierarchy via the gated mechanism.

\section{Conclusions and Discussions}

In this study, we introduced OSCAR, a novel deep learning framework for multi-omic cancer subtype classification that explicitly incorporates biological domain knowledge into its architectural design. By implementing an Ordered Cross-Attention (OCA) mechanism within OSCAR, we successfully modeled the directional regulatory flow dictated by the biological Central Dogma  $Genomics  \rightarrow Transcriptomics \rightarrow Proteomics$ . This hierarchical approach provides a powerful inductive bias, allowing the model to capture the complex, layer-wise dependencies between genomic alterations and their eventual proteomic functional outcomes, which are often neglected by symmetric or ``flat" integration strategies such as MOT \cite{MOT2023}.

Our experimental evaluation across three major TCGA cohorts—BRCA, LUNG, and COADREAD—demonstrates that the proposed framework consistently achieves SOTA performance. Specifically, our model reached a peak Accuracy of 91.27\% in the high-heterogeneity BRCA dataset and 96.71\% in LUNG subtype classification, outperforming established baselines such as MoGCN \cite{MoGCN2022}, MOGONET \cite{MOGONET2021}, and MoxGATE \cite{MoxGATE2023}. The integration of a Gated Fusion module was found to be critical, providing the flexibility to adapt to post-transcriptional noise and cases where the hierarchical coupling between layers is weakened.

The performance of deep learning models on multi-omic data is heavily influenced by class imbalance and the ``small $N$, large $P$" problem. As summarized in Table \ref{tab:dataset_summary}, the BRCA dataset exhibits extreme class imbalance, with the \textit{Luminal A} subtype representing nearly 48\% of the cohort while the \textit{HER2-enriched} subtype comprises only 10.4\%. This disproportion poses a significant challenge for symmetric attention models, which often converge toward the majority class. The superior Macro-F1 scores achieved by OSCAR demonstrate its robustness in identifying minority subtype signatures despite this skewed distribution.

Despite the architectural advantages of the Ordered Cross-Attention mechanism, extreme class imbalance—particularly in the BRCA dataset where the \textit{HER2-enriched} subtype represents only 10.4\% of the total samples—remains a hurdle. In our preliminary experiments, we attempted to replace the standard Cross-Entropy loss with the Focal Loss proposed by Lin et al.~\cite{FocalLoss2017} to encourage the model to focus on these hard-to-classify minority samples. The Focal Loss is defined as $FL(p_t) = -(1-p_t)^\gamma \log(p_t)$, where $\gamma$ is a focusing parameter. We tested values of $\gamma \in [0.5, 2, 5]$; however, while Focal Loss marginally improved the recall of the minority class, it led to a non-negligible drop in the precision of the majority classes (\textit{Luminal A} and \textit{Luminal B}). This suggests that in high-dimensional multi-omic spaces, simply re-weighting the loss function may not be sufficient to overcome the lack of sample diversity in minority cohorts.

To further enhance model robustness, future iterations of OSCAR will investigate alternative strategies such as: (1) \textit{Synthetic Data Generation} using GANs or VAEs to synthesize realistic profiles for under-represented subtypes; (2) \textit{Cost-Sensitive Attention} to adaptively increase query strength for minority classes; and (3) \textit{Transfer Learning} on larger, non-labeled pan-cancer datasets to learn more robust feature representations.

The high accuracy and stability of the OSCAR framework suggest that it could serve as a decision-support tool in precision oncology. By integrating genomic and proteomic signals, clinicians could better distinguish between subtypes that appear histologically similar but require different therapeutic strategies. Complementarily, our BRCA interpretability analysis shows that ordered multi-omic fusion yields clinically coherent saliency programs (luminal proliferation and the HER2 17q12 amplicon) and can nominate under-emphasized candidates such as \textit{TBC1D31} and \textit{GAS2L3} for hypothesis-driven follow-up. Future iterations of this work will aim to incorporate DNA methylation as an initial layer in the hierarchy, as epigenetic modifications often precede genomic structural changes. Additionally, we plan to extend OSCAR to Survival Analysis tasks, moving beyond classification to predict patient prognosis and treatment response durations.


% \section*{References}

% \bibliographystyle{apalike}
% \bibliography{reference}

\bibliographystyle{ieeetr} 
\bibliography{reference}


% \section{The Design, Intent, and \\ Limitations of the Templates}
% The templates are intended to {\bf{approximate the final look and page length of the articles/papers}}. {\bf{They are NOT intended to be the final produced work that is displayed in print or on IEEEXplore\textsuperscript{\textregistered}}}. They will help to give the authors an approximation of the number of pages that will be in the final version. The structure of the \LaTeX\ files, as designed, enable easy conversion to XML for the composition systems used by the IEEE. The XML files are used to produce the final print/IEEEXplore pdf and then converted to HTML for IEEEXplore.

% \section{Where to Get \LaTeX \ Help --- User Groups}
% The following online groups are helpful to beginning and experienced \LaTeX\ users. A search through their archives can provide many answers to common questions.
% \begin{list}{}{}
% \item{\url{http://www.latex-community.org/}} 
% \item{\url{https://tex.stackexchange.com/} }
% \end{list}

% \section{Other Resources}
% See \cite{ref1,ref2,ref3,ref4,ref5} for resources on formatting math into text and additional help in working with \LaTeX .

% \section{Text}
% For some of the remainer of this sample we will use dummy text to fill out paragraphs rather than use live text that may violate a copyright.

% Itam, que ipiti sum dem velit la sum et dionet quatibus apitet voloritet audam, qui aliciant voloreicid quaspe volorem ut maximusandit faccum conemporerum aut ellatur, nobis arcimus.
% Fugit odi ut pliquia incitium latum que cusapere perit molupta eaquaeria quod ut optatem poreiur? Quiaerr ovitior suntiant litio bearciur?

% Onseque sequaes rectur autate minullore nusae nestiberum, sum voluptatio. Et ratem sequiam quaspername nos rem repudandae volum consequis nos eium aut as molupta tectum ulparumquam ut maximillesti consequas quas inctia cum volectinusa porrum unt eius cusaest exeritatur? Nias es enist fugit pa vollum reium essusam nist et pa aceaqui quo elibusdandis deligendus que nullaci lloreri bla que sa coreriam explacc atiumquos simolorpore, non prehendunt lam que occum\cite{ref6} si aut aut maximus eliaeruntia dia sequiamenime natem sendae ipidemp orehend uciisi omnienetus most verum, ommolendi omnimus, est, veni aut ipsa volendelist mo conserum volores estisciis recessi nveles ut poressitatur sitiis ex endi diti volum dolupta aut aut odi as eatquo cullabo remquis toreptum et des accus dolende pores sequas dolores tinust quas expel moditae ne sum quiatis nis endipie nihilis etum fugiae audi dia quiasit quibus.
% \IEEEpubidadjcol
% Ibus el et quatemo luptatque doluptaest et pe volent rem ipidusa eribus utem venimolorae dera qui acea quam etur aceruptat.
% Gias anis doluptaspic tem et aliquis alique inctiuntiur?

% Sedigent, si aligend elibuscid ut et ium volo tem eictore pellore ritatus ut ut ullatus in con con pere nos ab ium di tem aliqui od magnit repta volectur suntio. Nam isquiante doluptis essit, ut eos suntionsecto debitiur sum ea ipitiis adipit, oditiore, a dolorerempos aut harum ius, atquat.

% Rum rem ditinti sciendunti volupiciendi sequiae nonsect oreniatur, volores sition ressimil inus solut ea volum harumqui to see\eqref{deqn_ex1a} mint aut quat eos explis ad quodi debis deliqui aspel earcius.

% \begin{equation}
% \label{deqn_ex1a}
% x = \sum_{i=0}^{n} 2{i} Q.
% \end{equation}

% Alis nime volorempera perferi sitio denim repudae pre ducilit atatet volecte ssimillorae dolore, ut pel ipsa nonsequiam in re nus maiost et que dolor sunt eturita tibusanis eatent a aut et dio blaudit reptibu scipitem liquia consequodi od unto ipsae. Et enitia vel et experferum quiat harum sa net faccae dolut voloria nem. Bus ut labo. Ita eum repraer rovitia samendit aut et volupta tecupti busant omni quiae porro que nossimodic temquis anto blacita conse nis am, que ereperum eumquam quaescil imenisci quae magnimos recus ilibeaque cum etum iliate prae parumquatemo blaceaquiam quundia dit apienditem rerit re eici quaes eos sinvers pelecabo. Namendignis as exerupit aut magnim ium illabor roratecte plic tem res apiscipsam et vernat untur a deliquaest que non cus eat ea dolupiducim fugiam volum hil ius dolo eaquis sitis aut landesto quo corerest et auditaquas ditae voloribus, qui optaspis exero cusa am, ut plibus.


% \section{Some Common Elements}
% \subsection{Sections and Subsections}
% Enumeration of section headings is desirable, but not required. When numbered, please be consistent throughout the article, that is, all headings and all levels of section headings in the article should be enumerated. Primary headings are designated with Roman numerals, secondary with capital letters, tertiary with Arabic numbers; and quaternary with lowercase letters. Reference and Acknowledgment headings are unlike all other section headings in text. They are never enumerated. They are simply primary headings without labels, regardless of whether the other headings in the article are enumerated. 

% \subsection{Citations to the Bibliography}
% The coding for the citations is made with the \LaTeX\ $\backslash${\tt{cite}} command. 
% This will display as: see \cite{ref1}.

% For multiple citations code as follows: {\tt{$\backslash$cite\{ref1,ref2,ref3\}}}
%  which will produce \cite{ref1,ref2,ref3}. For reference ranges that are not consecutive code as {\tt{$\backslash$cite\{ref1,ref2,ref3,ref9\}}} which will produce  \cite{ref1,ref2,ref3,ref9}

% \subsection{Lists}
% In this section, we will consider three types of lists: simple unnumbered, numbered, and bulleted. There have been many options added to IEEEtran to enhance the creation of lists. If your lists are more complex than those shown below, please refer to the original ``IEEEtran\_HOWTO.pdf'' for additional options.\\

% \subsubsection*{\bf A plain  unnumbered list}
% \begin{list}{}{}
% \item{bare\_jrnl.tex}
% \item{bare\_conf.tex}
% \item{bare\_jrnl\_compsoc.tex}
% \item{bare\_conf\_compsoc.tex}
% \item{bare\_jrnl\_comsoc.tex}
% \end{list}

% \subsubsection*{\bf A simple numbered list}
% \begin{enumerate}
% \item{bare\_jrnl.tex}
% \item{bare\_conf.tex}
% \item{bare\_jrnl\_compsoc.tex}
% \item{bare\_conf\_compsoc.tex}
% \item{bare\_jrnl\_comsoc.tex}
% \end{enumerate}

% \subsubsection*{\bf A simple bulleted list}
% \begin{itemize}
% \item{bare\_jrnl.tex}
% \item{bare\_conf.tex}
% \item{bare\_jrnl\_compsoc.tex}
% \item{bare\_conf\_compsoc.tex}
% \item{bare\_jrnl\_comsoc.tex}
% \end{itemize}





% \subsection{Figures}
% Fig. 1 is an example of a floating figure using the graphicx package.
%  Note that $\backslash${\tt{label}} must occur AFTER (or within) $\backslash${\tt{caption}}.
%  For figures, $\backslash${\tt{caption}} should occur after the $\backslash${\tt{includegraphics}}.

% \begin{figure}[!t]
% \centering
% \includegraphics[width=2.5in]{fig1}
% \caption{Simulation results for the network.}
% \label{fig_1}
% \end{figure}

% Fig. 2(a) and 2(b) is an example of a double column floating figure using two subfigures.
%  (The subfig.sty package must be loaded for this to work.)
%  The subfigure $\backslash${\tt{label}} commands are set within each subfloat command,
%  and the $\backslash${\tt{label}} for the overall figure must come after $\backslash${\tt{caption}}.
%  $\backslash${\tt{hfil}} is used as a separator to get equal spacing.
%  The combined width of all the parts of the figure should do not exceed the text width or a line break will occur.
% %
% \begin{figure*}[!t]
% \centering
% \subfloat[]{\includegraphics[width=2.5in]{fig1}%
% \label{fig_first_case}}
% \hfil
% \subfloat[]{\includegraphics[width=2.5in]{fig1}%
% \label{fig_second_case}}
% \caption{Dae. Ad quatur autat ut porepel itemoles dolor autem fuga. Bus quia con nessunti as remo di quatus non perum que nimus. (a) Case I. (b) Case II.}
% \label{fig_sim}
% \end{figure*}

% Note that often IEEE papers with multi-part figures do not place the labels within the image itself (using the optional argument to $\backslash${\tt{subfloat}}[]), but instead will
%  reference/describe all of them (a), (b), etc., within the main caption.
%  Be aware that for subfig.sty to generate the (a), (b), etc., subfigure
%  labels, the optional argument to $\backslash${\tt{subfloat}} must be present. If a
%  subcaption is not desired, leave its contents blank,
%  e.g.,$\backslash${\tt{subfloat}}[].


 

% \section{Tables}
% Note that, for IEEE-style tables, the
%  $\backslash${\tt{caption}} command should come BEFORE the table. Table captions use title case. Articles (a, an, the), coordinating conjunctions (and, but, for, or, nor), and most short prepositions are lowercase unless they are the first or last word. Table text will default to $\backslash${\tt{footnotesize}} as
%  the IEEE normally uses this smaller font for tables.
%  The $\backslash${\tt{label}} must come after $\backslash${\tt{caption}} as always.
 
% \begin{table}[!t]
% \caption{An Example of a Table\label{tab:table1}}
% \centering
% \begin{tabular}{|c||c|}
% \hline
% One & Two\\
% \hline
% Three & Four\\
% \hline
% \end{tabular}
% \end{table}

% \section{Algorithms}
% Algorithms should be numbered and include a short title. They are set off from the text with rules above and below the title and after the last line.

% \begin{algorithm}[H]
% \caption{Weighted Tanimoto ELM.}\label{alg:alg1}
% \begin{algorithmic}
% \STATE 
% \STATE {\textsc{TRAIN}}$(\mathbf{X} \mathbf{T})$
% \STATE \hspace{0.5cm}$ \textbf{select randomly } W \subset \mathbf{X}  $
% \STATE \hspace{0.5cm}$ N_\mathbf{t} \gets | \{ i : \mathbf{t}_i = \mathbf{t} \} | $ \textbf{ for } $ \mathbf{t}= -1,+1 $
% \STATE \hspace{0.5cm}$ B_i \gets \sqrt{ \textsc{max}(N_{-1},N_{+1}) / N_{\mathbf{t}_i} } $ \textbf{ for } $ i = 1,...,N $
% \STATE \hspace{0.5cm}$ \hat{\mathbf{H}} \gets  B \cdot (\mathbf{X}^T\textbf{W})/( \mathbb{1}\mathbf{X} + \mathbb{1}\textbf{W} - \mathbf{X}^T\textbf{W} ) $
% \STATE \hspace{0.5cm}$ \beta \gets \left ( I/C + \hat{\mathbf{H}}^T\hat{\mathbf{H}} \right )^{-1}(\hat{\mathbf{H}}^T B\cdot \mathbf{T})  $
% \STATE \hspace{0.5cm}\textbf{return}  $\textbf{W},  \beta $
% \STATE 
% \STATE {\textsc{PREDICT}}$(\mathbf{X} )$
% \STATE \hspace{0.5cm}$ \mathbf{H} \gets  (\mathbf{X}^T\textbf{W} )/( \mathbb{1}\mathbf{X}  + \mathbb{1}\textbf{W}- \mathbf{X}^T\textbf{W}  ) $
% \STATE \hspace{0.5cm}\textbf{return}  $\textsc{sign}( \mathbf{H} \beta )$
% \end{algorithmic}
% \label{alg1}
% \end{algorithm}

% Que sunt eum lam eos si dic to estist, culluptium quid qui nestrum nobis reiumquiatur minimus minctem. Ro moluptat fuga. Itatquiam ut laborpo rersped exceres vollandi repudaerem. Ulparci sunt, qui doluptaquis sumquia ndestiu sapient iorepella sunti veribus. Ro moluptat fuga. Itatquiam ut laborpo rersped exceres vollandi repudaerem. 
% \section{Mathematical Typography \\ and Why It Matters}

% Typographical conventions for mathematical formulas have been developed to {\bf provide uniformity and clarity of presentation across mathematical texts}. This enables the readers of those texts to both understand the author's ideas and to grasp new concepts quickly. While software such as \LaTeX \ and MathType\textsuperscript{\textregistered} can produce aesthetically pleasing math when used properly, it is also very easy to misuse the software, potentially resulting in incorrect math display.

% IEEE aims to provide authors with the proper guidance on mathematical typesetting style and assist them in writing the best possible article. As such, IEEE has assembled a set of examples of good and bad mathematical typesetting \cite{ref1,ref2,ref3,ref4,ref5}. 

% Further examples can be found at \url{http://journals.ieeeauthorcenter.ieee.org/wp-content/uploads/sites/7/IEEE-Math-Typesetting-Guide-for-LaTeX-Users.pdf}

% \subsection{Display Equations}
% The simple display equation example shown below uses the ``equation'' environment. To number the equations, use the $\backslash${\tt{label}} macro to create an identifier for the equation. LaTeX will automatically number the equation for you.
% \begin{equation}
% \label{deqn_ex1}
% x = \sum_{i=0}^{n} 2{i} Q.
% \end{equation}

% \noindent is coded as follows:
% \begin{verbatim}
% \begin{equation}
% \label{deqn_ex1}
% x = \sum_{i=0}^{n} 2{i} Q.
% \end{equation}
% \end{verbatim}

% To reference this equation in the text use the $\backslash${\tt{ref}} macro. 
% Please see (\ref{deqn_ex1})\\
% \noindent is coded as follows:
% \begin{verbatim}
% Please see (\ref{deqn_ex1})\end{verbatim}

% \subsection{Equation Numbering}
% {\bf{Consecutive Numbering:}} Equations within an article are numbered consecutively from the beginning of the
% article to the end, i.e., (1), (2), (3), (4), (5), etc. Do not use roman numerals or section numbers for equation numbering.

% \noindent {\bf{Appendix Equations:}} The continuation of consecutively numbered equations is best in the Appendix, but numbering
%  as (A1), (A2), etc., is permissible.\\

% \noindent {\bf{Hyphens and Periods}}: Hyphens and periods should not be used in equation numbers, i.e., use (1a) rather than
% (1-a) and (2a) rather than (2.a) for subequations. This should be consistent throughout the article.

% \subsection{Multi-Line Equations and Alignment}
% Here we show several examples of multi-line equations and proper alignments.

% \noindent {\bf{A single equation that must break over multiple lines due to length with no specific alignment.}}
% \begin{multline}
% \text{The first line of this example}\\
% \text{The second line of this example}\\
% \text{The third line of this example}
% \end{multline}

% \noindent is coded as:
% \begin{verbatim}
% \begin{multline}
% \text{The first line of this example}\\
% \text{The second line of this example}\\
% \text{The third line of this example}
% \end{multline}
% \end{verbatim}

% \noindent {\bf{A single equation with multiple lines aligned at the = signs}}
% \begin{align}
% a &= c+d \\
% b &= e+f
% \end{align}
% \noindent is coded as:
% \begin{verbatim}
% \begin{align}
% a &= c+d \\
% b &= e+f
% \end{align}
% \end{verbatim}

% The {\tt{align}} environment can align on multiple  points as shown in the following example:
% \begin{align}
% x &= y & X & =Y & a &=bc\\
% x' &= y' & X' &=Y' &a' &=bz
% \end{align}
% \noindent is coded as:
% \begin{verbatim}
% \begin{align}
% x &= y & X & =Y & a &=bc\\
% x' &= y' & X' &=Y' &a' &=bz
% \end{align}
% \end{verbatim}





% \subsection{Subnumbering}
% The amsmath package provides a {\tt{subequations}} environment to facilitate subnumbering. An example:

% \begin{subequations}\label{eq:2}
% \begin{align}
% f&=g \label{eq:2A}\\
% f' &=g' \label{eq:2B}\\
% \mathcal{L}f &= \mathcal{L}g \label{eq:2c}
% \end{align}
% \end{subequations}

% \noindent is coded as:
% \begin{verbatim}
% \begin{subequations}\label{eq:2}
% \begin{align}
% f&=g \label{eq:2A}\\
% f' &=g' \label{eq:2B}\\
% \mathcal{L}f &= \mathcal{L}g \label{eq:2c}
% \end{align}
% \end{subequations}

% \end{verbatim}

% \subsection{Matrices}
% There are several useful matrix environments that can save you some keystrokes. See the example coding below and the output.

% \noindent {\bf{A simple matrix:}}
% \begin{equation}
% \begin{matrix}  0 &  1 \\ 
% 1 &  0 \end{matrix}
% \end{equation}
% is coded as:
% \begin{verbatim}
% \begin{equation}
% \begin{matrix}  0 &  1 \\ 
% 1 &  0 \end{matrix}
% \end{equation}
% \end{verbatim}

% \noindent {\bf{A matrix with parenthesis}}
% \begin{equation}
% \begin{pmatrix} 0 & -i \\
%  i &  0 \end{pmatrix}
% \end{equation}
% is coded as:
% \begin{verbatim}
% \begin{equation}
% \begin{pmatrix} 0 & -i \\
%  i &  0 \end{pmatrix}
% \end{equation}
% \end{verbatim}

% \noindent {\bf{A matrix with square brackets}}
% \begin{equation}
% \begin{bmatrix} 0 & -1 \\ 
% 1 &  0 \end{bmatrix}
% \end{equation}
% is coded as:
% \begin{verbatim}
% \begin{equation}
% \begin{bmatrix} 0 & -1 \\ 
% 1 &  0 \end{bmatrix}
% \end{equation}
% \end{verbatim}

% \noindent {\bf{A matrix with curly braces}}
% \begin{equation}
% \begin{Bmatrix} 1 &  0 \\ 
% 0 & -1 \end{Bmatrix}
% \end{equation}
% is coded as:
% \begin{verbatim}
% \begin{equation}
% \begin{Bmatrix} 1 &  0 \\ 
% 0 & -1 \end{Bmatrix}
% \end{equation}\end{verbatim}

% \noindent {\bf{A matrix with single verticals}}
% \begin{equation}
% \begin{vmatrix} a &  b \\ 
% c &  d \end{vmatrix}
% \end{equation}
% is coded as:
% \begin{verbatim}
% \begin{equation}
% \begin{vmatrix} a &  b \\ 
% c &  d \end{vmatrix}
% \end{equation}\end{verbatim}

% \noindent {\bf{A matrix with double verticals}}
% \begin{equation}
% \begin{Vmatrix} i &  0 \\ 
% 0 & -i \end{Vmatrix}
% \end{equation}
% is coded as:
% \begin{verbatim}
% \begin{equation}
% \begin{Vmatrix} i &  0 \\ 
% 0 & -i \end{Vmatrix}
% \end{equation}\end{verbatim}

% \subsection{Arrays}
% The {\tt{array}} environment allows you some options for matrix-like equations. You will have to manually key the fences, but there are other options for alignment of the columns and for setting horizontal and vertical rules. The argument to {\tt{array}} controls alignment and placement of vertical rules.

% A simple array
% \begin{equation}
% \left(
% \begin{array}{cccc}
% a+b+c & uv & x-y & 27\\
% a+b & u+v & z & 134
% \end{array}\right)
% \end{equation}
% is coded as:
% \begin{verbatim}
% \begin{equation}
% \left(
% \begin{array}{cccc}
% a+b+c & uv & x-y & 27\\
% a+b & u+v & z & 134
% \end{array} \right)
% \end{equation}
% \end{verbatim}

% A slight variation on this to better align the numbers in the last column
% \begin{equation}
% \left(
% \begin{array}{cccr}
% a+b+c & uv & x-y & 27\\
% a+b & u+v & z & 134
% \end{array}\right)
% \end{equation}
% is coded as:
% \begin{verbatim}
% \begin{equation}
% \left(
% \begin{array}{cccr}
% a+b+c & uv & x-y & 27\\
% a+b & u+v & z & 134
% \end{array} \right)
% \end{equation}
% \end{verbatim}

% An array with vertical and horizontal rules
% \begin{equation}
% \left( \begin{array}{c|c|c|r}
% a+b+c & uv & x-y & 27\\ \hline
% a+b & u+v & z & 134
% \end{array}\right)
% \end{equation}
% is coded as:
% \begin{verbatim}
% \begin{equation}
% \left(
% \begin{array}{c|c|c|r}
% a+b+c & uv & x-y & 27\\
% a+b & u+v & z & 134
% \end{array} \right)
% \end{equation}
% \end{verbatim}
% Note the argument now has the pipe "$\vert$" included to indicate the placement of the vertical rules.


% \subsection{Cases Structures}
% Many times cases can be miscoded using the wrong environment, i.e., {\tt{array}}. Using the {\tt{cases}} environment will save keystrokes (from not having to type the $\backslash${\tt{left}}$\backslash${\tt{lbrace}}) and automatically provide the correct column alignment.
% \begin{equation*}
% {z_m(t)} = \begin{cases}
% 1,&{\text{if}}\ {\beta }_m(t) \\ 
% {0,}&{\text{otherwise.}} 
% \end{cases}
% \end{equation*}
% \noindent is coded as follows:
% \begin{verbatim}
% \begin{equation*}
% {z_m(t)} = 
% \begin{cases}
% 1,&{\text{if}}\ {\beta }_m(t),\\ 
% {0,}&{\text{otherwise.}} 
% \end{cases}
% \end{equation*}
% \end{verbatim}
% \noindent Note that the ``\&'' is used to mark the tabular alignment. This is important to get  proper column alignment. Do not use $\backslash${\tt{quad}} or other fixed spaces to try and align the columns. Also, note the use of the $\backslash${\tt{text}} macro for text elements such as ``if'' and ``otherwise.''

% \subsection{Function Formatting in Equations}
% Often, there is an easy way to properly format most common functions. Use of the $\backslash$ in front of the function name will in most cases, provide the correct formatting. When this does not work, the following example provides a solution using the $\backslash${\tt{text}} macro:

% \begin{equation*} 
%   d_{R}^{KM} = \underset {d_{l}^{KM}} {\text{arg min}} \{ d_{1}^{KM},\ldots,d_{6}^{KM}\}.
% \end{equation*}

% \noindent is coded as follows:
% \begin{verbatim}
% \begin{equation*} 
%  d_{R}^{KM} = \underset {d_{l}^{KM}} 
%  {\text{arg min}} \{ d_{1}^{KM},
%  \ldots,d_{6}^{KM}\}.
% \end{equation*}
% \end{verbatim}

% \subsection{ Text Acronyms Inside Equations}
% This example shows where the acronym ``MSE" is coded using $\backslash${\tt{text\{\}}} to match how it appears in the text.

% \begin{equation*}
%  \text{MSE} = \frac {1}{n}\sum _{i=1}^{n}(Y_{i} - \hat {Y_{i}})^{2}
% \end{equation*}

% \begin{verbatim}
% \begin{equation*}
%  \text{MSE} = \frac {1}{n}\sum _{i=1}^{n}
% (Y_{i} - \hat {Y_{i}})^{2}
% \end{equation*}
% \end{verbatim}

% \section{Conclusion}
% The conclusion goes here.


% \section*{Acknowledgments}
% This should be a simple paragraph before the References to thank those individuals and institutions who have supported your work on this article.



% {\appendix[Proof of the Zonklar Equations]
% Use $\backslash${\tt{appendix}} if you have a single appendix:
% Do not use $\backslash${\tt{section}} anymore after $\backslash${\tt{appendix}}, only $\backslash${\tt{section*}}.
% If you have multiple appendixes use $\backslash${\tt{appendices}} then use $\backslash${\tt{section}} to start each appendix.
% You must declare a $\backslash${\tt{section}} before using any $\backslash${\tt{subsection}} or using $\backslash${\tt{label}} ($\backslash${\tt{appendices}} by itself
%  starts a section numbered zero.)}



% %{\appendices
% %\section*{Proof of the First Zonklar Equation}
% %Appendix one text goes here.
% % You can choose not to have a title for an appendix if you want by leaving the argument blank
% %\section*{Proof of the Second Zonklar Equation}
% %Appendix two text goes here.}



% \section{References Section}
% You can use a bibliography generated by BibTeX as a .bbl file.
%  BibTeX documentation can be easily obtained at:
%  http://mirror.ctan.org/biblio/bibtex/contrib/doc/
%  The IEEEtran BibTeX style support page is:
%  http://www.michaelshell.org/tex/ieeetran/bibtex/
 
%  % argument is your BibTeX string definitions and bibliography database(s)
% %\bibliography{IEEEabrv,../bib/paper}
% %
% \section{Simple References}
% You can manually copy in the resultant .bbl file and set second argument of $\backslash${\tt{begin}} to the number of references
%  (used to reserve space for the reference number labels box).

% \begin{thebibliography}{1}
% \bibliographystyle{IEEEtran}

% \bibitem{ref1}
% {\it{Mathematics Into Type}}. American Mathematical Society. [Online]. Available: https://www.ams.org/arc/styleguide/mit-2.pdf

% \bibitem{ref2}
% T. W. Chaundy, P. R. Barrett and C. Batey, {\it{The Printing of Mathematics}}. London, U.K., Oxford Univ. Press, 1954.

% \bibitem{ref3}
% F. Mittelbach and M. Goossens, {\it{The \LaTeX Companion}}, 2nd ed. Boston, MA, USA: Pearson, 2004.

% \bibitem{ref4}
% G. Gr\"atzer, {\it{More Math Into LaTeX}}, New York, NY, USA: Springer, 2007.

% \bibitem{ref5}M. Letourneau and J. W. Sharp, {\it{AMS-StyleGuide-online.pdf,}} American Mathematical Society, Providence, RI, USA, [Online]. Available: http://www.ams.org/arc/styleguide/index.html

% \bibitem{ref6}
% H. Sira-Ramirez, ``On the sliding mode control of nonlinear systems,'' \textit{Syst. Control Lett.}, vol. 19, pp. 303--312, 1992.

% \bibitem{ref7}
% A. Levant, ``Exact differentiation of signals with unbounded higher derivatives,''  in \textit{Proc. 45th IEEE Conf. Decis.
% Control}, San Diego, CA, USA, 2006, pp. 5585--5590. DOI: 10.1109/CDC.2006.377165.

% \bibitem{ref8}
% M. Fliess, C. Join, and H. Sira-Ramirez, ``Non-linear estimation is easy,'' \textit{Int. J. Model., Ident. Control}, vol. 4, no. 1, pp. 12--27, 2008.

% \bibitem{ref9}
% R. Ortega, A. Astolfi, G. Bastin, and H. Rodriguez, ``Stabilization of food-chain systems using a port-controlled Hamiltonian description,'' in \textit{Proc. Amer. Control Conf.}, Chicago, IL, USA,
% 2000, pp. 2245--2249.

% \end{thebibliography}


% \newpage

% \section{Biography Section}
% If you have an EPS/PDF photo (graphicx package needed), extra braces are
%  needed around the contents of the optional argument to biography to prevent
%  the LaTeX parser from getting confused when it sees the complicated
%  $\backslash${\tt{includegraphics}} command within an optional argument. (You can create
%  your own custom macro containing the $\backslash${\tt{includegraphics}} command to make things
%  simpler here.)
 
% \vspace{11pt}

% \bf{If you include a photo:}\vspace{-33pt}
% \begin{IEEEbiography}[{\includegraphics[width=1in,height=1.25in,clip,keepaspectratio]{fig1}}]{Michael Shell}
% Use $\backslash${\tt{begin\{IEEEbiography\}}} and then for the 1st argument use $\backslash${\tt{includegraphics}} to declare and link the author photo.
% Use the author name as the 3rd argument followed by the biography text.
% \end{IEEEbiography}

% \vspace{11pt}

% \bf{If you will not include a photo:}\vspace{-33pt}
% \begin{IEEEbiographynophoto}{John Doe}
% Use $\backslash${\tt{begin\{IEEEbiographynophoto\}}} and the author name as the argument followed by the biography text.
% \end{IEEEbiographynophoto}




% \vfill

\end{document}

